\chapter{E3: LSTM và GRU trên chuỗi ảnh}
\section{Biểu diễn và baseline}
Rows và columns có T=28,D=28; patches4$\times$4 có T=49,D=16. Các ảnh giữ cùng
split Fashion-MNIST. Notebook định nghĩa và train MLP/CNN-B riêng, không đọc
checkpoint E1. Điều này cho phép chạy từ kernel mới chỉ với dataset local đã xử lý.
\section{Recurrent cells và kiểm chứng}
LSTM dùng gates input/forget/output, candidate, hidden và cell state:
\begin{equation}c_t=f_t\odot c_{t-1}+i_t\odot g_t,\quad h_t=o_t\odot\tanh(c_t).\end{equation}
GRU dùng reset/update với convention PyTorch:
\begin{equation}n_t=\tanh(W_{in}x_t+b_{in}+r_t\odot(W_{hn}h_{t-1}+b_{hn})),\quad
h_t=(1-z_t)\odot n_t+z_t\odot h_{t-1}.\end{equation}
Điểm cần chú ý: reset nhân sau projection hidden, gồm hidden bias. Manual cell được
copy weight sang nn.LSTMCell/nn.GRUCell để so forward và gradient tất cả tham số.
Model dùng hidden128, một layer, hidden cuối và linear10.
\section{Protocol và ma trận}
LSTM/GRU$\times$ba representation tạo6 run; MLP/CNN-B thêm2. AdamW$10^{-3}$,
WD$10^{-4}$, batch128, tối đa 100 epoch, patience 4, gradient clipping1 cho recurrent.
Baseline giữ protocol riêng tương ứng kiến trúc, không clipping như recurrent.
Mỗi lần Run All tạo run mới, train trên toàn bộ train54k/val6k/test10k bằng CUDA.
Early stopping theo validation loss, min\_delta=$10^{-4}$, patience 4;
best checkpoint có validation loss thấp nhất. Text progress cập nhật tại chỗ,
xác nhận tensor/gradient CUDA, tiến độ batch, VRAM và early-stop counter.
Không dùng test cho lựa chọn. Log chi tiết, history, config và seed được lưu.
\section{Kết quả và phân tích}
\ResultTable{e3}{So sánh recurrent representation và baseline nội bộ.}
\FigureOrPlaceholder{assets/e3/sequence_representations.png}{Rows, columns, patches}{Minh họa thứ tự và feature của chuỗi ảnh.}
\FigureOrPlaceholder{assets/e3/comparison.png}{So sánh E3}{Accuracy validation của tám full run E3.}
\FigureOrPlaceholder{assets/e3/learning_curves_loss.png}{Learning curves loss}{Ba nhóm tokenizer/representation, màu theo model; nét đứt train, nét liền validation.}
\FigureOrPlaceholder{assets/e3/learning_curves_accuracy.png}{Learning curves accuracy}{Accuracy theo epoch của lần train mới.}
\FigureOrPlaceholder{assets/e3/recurrent_representation_comparison.png}{LSTM và GRU}{So sánh cùng representation: accuracy, capacity, giây/epoch.}
\FigureOrPlaceholder{assets/e3/baseline_curves.png}{So sánh baseline}{LSTM/GRU đại diện chọn bằng validation và MLP/CNN-B train tại chỗ.}
\FigureOrPlaceholder{assets/e3/confusion_test_pair_1.png}{Confusion matrix test, cặp 1}{Mỗi ô ghi số ảnh và tỷ lệ theo hàng; cùng thang 0--100\%.}
\FigureOrPlaceholder{assets/e3/confusion_test_pair_2.png}{Confusion matrix test, cặp 2}{Mỗi ô ghi số ảnh và tỷ lệ theo hàng; cùng thang 0--100\%.}
\FigureOrPlaceholder{assets/e3/confusion_test_pair_3.png}{Confusion matrix test, cặp 3}{Mỗi ô ghi số ảnh và tỷ lệ theo hàng; cùng thang 0--100\%.}
\FigureOrPlaceholder{assets/e3/confusion_test_pair_4.png}{Confusion matrix test, cặp 4}{Mỗi ô ghi số ảnh và tỷ lệ theo hàng; cùng thang 0--100\%.}
\FigureOrPlaceholder{assets/e3/failure_examples_test.png}{Ảnh lỗi theo model/lớp}{Model đại diện chọn bằng validation; mỗi lớp/model lấy lỗi đầu tiên.}
\Pending{Bổ sung kết quả thật, hiệu ứng thứ tự quét, loss gap và lỗi. So sánh chi phí
LSTM/GRU, capacity và locality với CNN; hidden state không phải giải thích nhân quả.}

\IfFileExists{assets/e3/test_results.tex}{\begin{table}[H]\centering\caption{Metric test sau khi khóa protocol.}\resizebox{\linewidth}{!}{\input{assets/e3/test_results.tex}}\end{table}}{}
